% Direct-authority diplomatic French for printed p.188.
\clearpage
\SourcePage{188}

\LayerTarget{tohoku:remark:ext-bifunctoriality:p1}
\noindent les $\mathrm{Ext}_{\mathbf O}^p(A,B)$ forment non seulement des
\textit{foncteurs cohomologiques covariants par rapport à $B$}, mais aussi des
\textit{foncteurs cohomologiques contravariants par rapport à $A$}, grâce au
fait que pour $B$ injectif, les foncteurs
$A\to\mathrm{Hom}_{\mathbf O}(A,B)$ et
$A\to\mathbf{Hom}_{\mathbf O}(A,B)$ sont exacts. (Le deuxième fait se vérifie
immédiatement à partir du premier, grâce à prop. 3.1.3).

\LayerTarget{tohoku:remark:ext-center-modules}
Remarquons enfin, pour finir les sorites, que les
$\mathrm{Ext}_{\mathbf O}^p(X;A,B)$ sont des modules sur le centre de $\Gamma({\mathbf O})$, et
que les $\mathbf{Ext}_{\mathbf O}^p(A,B)$ sont des Modules sur le faisceau d'anneaux
${\mathbf O}^\natural$ centre de ${\mathbf O}$.

\LayerTarget{tohoku:construction:hom-composite-functors}
La formule (4.1.2) peut aussi s'écrire, en désignant par $h_A$ resp.
$\mathbf{h}_A$ le foncteur covariant de $\mathbf C^{\mathbf O}$ dans la catégorie $\mathbf C$ des
groupes abéliens, resp. $\mathbf C^X$, défini par
\LayerTarget{tohoku:eq:4.2.4}
\[
\tag{42.4}
\begin{gathered}
h_A(B)=\mathrm{Hom}_{\mathbf O}(A,B),\qquad
\mathbf{h}_A(B)=\mathbf{Hom}_{\mathbf O}(A,B):\\
h_A=\Gamma\mathbf{h}_A.
\end{gathered}
\]
Donc $h_A$ apparaît comme un foncteur composé, qui est justiciable de
théorème 2.4.1 en vertu de prop. 4.1.3 et du corollaire à la prop. 3.3.1. On
obtient donc :

\LayerTarget{tohoku:theorem:4.2.1}
\textsc{Théorème 4.2.1.} \textit{Soit $X$ un espace topologique muni
d'un faisceau d'anneaux avec unité ${\mathbf O}$, soit $\mathbf C^{\mathbf O}$ la catégorie des
${\mathbf O}$-Modules sur $X$. Soit $A\in \mathbf C^{\mathbf O}$ un ${\mathbf O}$-module fixé. Alors il existe sur
$\mathbf C^{\mathbf O}$ un foncteur spectral cohomologique, aboutissant au foncteur gradué
$(\mathrm{Ext}_{\mathbf O}^p(X;A,B))$, et dont le terme initial est}
\LayerTarget{tohoku:eq:4.2.5}
\[
\tag{4.2.5}
I_2^{p,q}(A,B)=H^p(X,\mathbf{Ext}_{\mathbf O}^q(A,B)).
\]

\LayerTarget{tohoku:eq:4.2.6}
On en conclut en particulier des ``\textit{edge homomorphisms}''
\[
\tag{4.2.6}
\begin{aligned}
H^p(X,\mathbf{Hom}_{\mathbf O}(A,B))&\to
  \mathrm{Ext}_{\mathbf O}^p(X;A,B),\\
\mathrm{Ext}_{\mathbf O}^p(X;A,B)&\to
  H^0(X,\mathbf{Ext}_{\mathbf O}^p(A,B))
\end{aligned}
\]
(ce dernier étant d'ailleurs celui qui résulte de l'interprétation de
$\mathbf{Ext}_{\mathbf O}^p(A,B)$ comme le faisceau associé au préfaisceau
$\mathrm{Ext}_{\mathbf O}^p(-;A,B)$), et une suite exacte à 5 termes :
\LayerTarget{tohoku:eq:4.2.7}
\[
\tag{4.2.7}
\begin{aligned}
0\to H^1(X,\mathbf{Hom}_{\mathbf O}(A,B))
&\to \mathrm{Ext}_{\mathbf O}^1(X;A,B)
\to H^0(X,\mathbf{Ext}_{\mathbf O}^1(A,B))\\
&\to H^2(X,\mathbf{Hom}_{\mathbf O}(A,B))
\to \mathrm{Ext}_{\mathbf O}^2(X;A,B).
\end{aligned}
\]
\LayerTarget{tohoku:remark:ext1-extension-classes}
Cette suite exacte \textit{précise en particulier la structure du groupe
$\mathrm{Ext}_{\mathbf O}^1(X;A,B)$ des classes de ${\mathbf O}$-Modules extensions de $A$
\textup{(Module quotient)} par $B$ \textup{(sous-Module).}}

\LayerTarget{tohoku:construction:stalk-ext-comparison}
Pour pouvoir utiliser la suite spectrale obtenue, il faut préciser encore le
calcul des $\mathbf{Ext}_{\mathbf O}^p(A,B)$. Comme le foncteur $F\to F(x)$ qui,
à un faisceau abélien $F$, associe son groupe ponctuel $F(x)$, est \textit{exact}, les
$\mathbf{Ext}_{\mathbf O}^p(A,B)(x)$ peuvent être considérés, pour $A$ fixé et $p$, $B$
variables, comme formant un foncteur cohomologique covariant universel en $B$,
se réduisant à $\mathbf{Hom}_{\mathbf O}(A,B)(x)$ en dimension 0. De même, les
$\mathrm{Ext}_{{\mathbf O}(x)}^p(A(x),B(x))$ forment un foncteur cohomologique en $B$,
se réduisant à $\mathrm{Hom}_{{\mathbf O}(x)}(A(x),B(x))$ en dimension 0. On conclut
donc, du homomorphisme fonctoriel (4.1.3), (grâce à la définition des
foncteurs cohomologiques universels) des homomorphismes
\LayerTarget{tohoku:eq:4.2.8}
\[
\tag{4.2.8}
\mathrm{Ext}_{\mathbf O}^p(A,B)(x)\to
\mathrm{Ext}_{{\mathbf O}(x)}^p(A(x),B(x))
\]

\TohokuContinuation{189}
